\documentclass[10pt,a4paper]{amsart}
\usepackage[margin=19mm]{geometry}
\usepackage{amsmath,amssymb,mathtools}
\usepackage[scr=rsfs,cal=euler]{mathalfa}
\usepackage{xfrac}
\usepackage[unicode,hidelinks,bookmarks=false]{hyperref}
\newcommand{\ie}{i.e.\ }
\newcommand{\cf}{cf.\ }
\newcommand{\resp}{respectively\ }
\newcommand{\Subsection}{\subsection}
\providecommand{\nobreakdash}{\nobreak\hbox{-}}
\setlength{\emergencystretch}{2em}
\multlinegap0pt
\usepackage{amssymb}
\usepackage{mathtools}
\usepackage{xcolor}
\newtheorem{lem}{Lemma}
\newcommand{\cN}{\mathcal{N}}
\newcommand{\fg}{\mathfrak{g}}
\title{Counting points on generic character varieties}
\author{Masoud Kamgarpour, GyeongHyeon Nam, Bailey Whitbread,, Stefano Giannini}
\date{Source: arXiv:2409.04735v4 (CC BY 4.0)}
\begin{document}
\maketitle

\noindent\emph{Source and licence.} Counting points on generic character varieties. Version arXiv:2409.04735v4 (CC BY 4.0), distributed by arXiv under the Creative Commons Attribution 4.0 International licence, http://creativecommons.org/licenses/by/4.0/, as recorded in the arXiv licence metadata for this exact version and shown on the abstract page https://arxiv.org/abs/2409.04735v4. Full licence notice: \texttt{licenses/arxiv-2409.04735-CC-BY-4.0.txt}.\par\smallskip

\noindent\textit{Editorial scope.} Counted unit: the article's lem lem (source0.tex lines 1423--1425). The article's complete original proof of it is delivered with it (source0.tex lines 1426--1436, one \texttt{proof} environment of 11 lines). No mathematics is written, repaired or re-derived: the statement and the proof are the article's own lines, in the article's own notation, and no retained line is substituted or re-flowed.  No further source-local context is needed: every cross-reference of every retained line resolves inside the two delivered spans.  Nothing was omitted from any retained span, so the package carries no omission record.  No retained line contains a citation, so the wrapper carries no bibliography and no bibliography entry.  No retained line contains a figure environment, an image-inclusion command or drawing code, so no image is delivered and the package declares no asset dependency.  Editorial formatting only, in addition: an AMS \texttt{amsart} preamble that restates the article's own declarations and macros used by the retained lines, namely the environments \texttt{lem} on separate counters, together with the standard packages those lines need.  The retained environments print in this document as Lemma 1 in order of appearance, whereas the article prints the same items with the numbering of its own full document.  The complete native source of the article as distributed in the arXiv e-print for this exact version is bundled unchanged as \texttt{sources/164914/source0.tex}.
\begin{lem} If $\mathfrak{O}$ has type $\tau=[L,\cN]$, then 
$\displaystyle |\mathfrak{O}| =\frac{|G^F||\cN|}{|L^F|}$. 
\end{lem} 

\begin{proof} 
Let $x=x_s+x_n\in \fg^F$.  Jordan decomposition is preserved under adjoint action; thus, for $g\in G$, $g\cdot x=g\cdot x_s+g\cdot x_n$ is the Jordan decomposition of $g\cdot x$. It follows that  $g\cdot x=x$ if and only if $g\cdot x_s=x_s$ and $g\cdot x_n=x_n$. In other words,  
\[
G_x=G_{x_s}\cap G_{x_n} = C_{G_{x_s}}(x_n).
\]
 Now if $x\in \fg^F$ has type $\tau=[L,\cN]$, then 
$C_{G_{x_s}}(x_n)=|L^F|/|\cN|$. Thus, if $\mathfrak{O}$ is the $G^F$-orbit of $x$, we conclude 
\[
|\mathfrak{O}| = \frac{|G^F|}{|G_{x}^F|}=\frac{|G^F|}{|C_{G_{x_s}}(x_n)^F|}=\frac{|G^F||\cN|}{|L^F|}. \qedhere
\]
\end{proof} 

\end{document}
